%
% dia3.tex
% 
% Copyright 2014 Rony J. Letona <rony@zronyj.com>
% 
% This program is free software; you can redistribute it and/or modify
% it under the terms of the GNU General Public License as published by
% the Free Software Foundation; either version 2 of the License, or
% any later version.
% 
% This program is distributed in the hope that it will be useful,
% but WITHOUT ANY WARRANTY; without even the implied warranty of
% MERCHANTABILITY or FITNESS FOR A PARTICULAR PURPOSE.  See the
% GNU General Public License for more details.
% 
% You should have received a copy of the GNU General Public License
% along with this program; if not, write to the Free Software
% Foundation, Inc., 51 Franklin Street, Fifth Floor, Boston,
% MA 02110-1301, USA.
%

\documentclass[10pt,letterpaper]{article}
\usepackage[latin1]{inputenc}
\usepackage[spanish]{babel}
\usepackage{graphicx}
\usepackage{hyperref}
\usepackage{amsmath}
\usepackage{amsfonts}
\usepackage{amssymb}
\usepackage{color}
\usepackage{float}
\usepackage[left=2cm,right=2cm,top=2cm,bottom=2cm]{geometry}
\author{Rony J. Letona}
\title{Taller de Qu\'imica Computacional Aplicada: D\'ia 3}
\definecolor{light-gray}{gray}{0.90}

\newcommand{\tab}[1]{\hspace{.2\textwidth}\rlap{#1}}

\newcommand{\inlinecode}[1]{
\colorbox{light-gray}{\texttt{#1}}
}

\newsavebox{\selvestebox}
\newenvironment{Code}
{
\begin{lrbox}{\selvestebox}%
\begin{minipage}{\dimexpr\columnwidth-2\fboxsep\relax}
\fontfamily{\ttdefault}\selectfont
}
{\end{minipage}\end{lrbox}%
\begin{center}
\colorbox{light-gray}{\usebox{\selvestebox}}
\end{center}
}

\newcommand{\Picture}[1]
{
	\begin{figure}[H]
	\begin{flushleft}
	\includegraphics[width=\columnwidth]{#1}
	\end{flushleft}
	\end{figure}
}

\begin{document}
\maketitle

\section{Ejercicios con el Sistema de Control de Revisi\'on - Git}
Al comenzar a trabajar en un proyecto en computaci\'on, generalmente se trabaja con archivos sencillos de texto, como hemos visto anteriormente. Estos tienen el c\'odigo que escribimos y al compilarlos/interpretarlos y ejecutarlos, resultan en programas. Todo funciona as\'i. A veces el c\'odigo en los archivos es comprensible, a veces no lo es porque no es para que lo entendamos nosotros sino el ordenador. Sin embargo, a todos nos ha pasado que borramos un archivo importante o que cambiamos algo que no deb\'iamos de cambiar y no hay manera de recuperarlo. Desde documentos como reportes de laboratorio hasta cartas, siempre nos pasa en alguna ocasi\'on que perdemos informaci\'on importante. Para evitar esto en ambientes de desarrollo, se dise\~n\'o una alternativa.\\

Git es la soluci\'on m\'as eficiente. Es un sistema en el que cada revisi\'on o versi\'on de un archivo se va guardando gradualmente cuando nosotros lo indiquemos; se lleva registro de cu\'ales cambios se hicieron y de cu\'ando se hicieron. Es un cuaderno de laboratorio digital! No nos salvar\'a un programa a medio escribir si perdemos acceso a nuestros documentos o si deseamos trabajar en colaboraci\'on con m\'as personas. Es un sistema para llevar autom\'aticamente registro del trabajo realizado. Lo usan grandes empresas, as\'i como peque\~nas iniciativas en sus proyectos. Y no se tiene que limitar a software, puede tratarse de libros, art\'iculos, bases de datos peque\~nas y cualquier otra cosa que sufra cambios con el tiempo o requiera ser compartida. Por ello, y por otras razones que veremos m\'as adelante, vamos a hacer algunos ejercicios con el sistema git.

\section{Creando un repositorio}
Lo primero que haremos ser\'a crear un repositorio. Esto \'ultimo es un lugar en donde vamos a colocar todos los archivos y documentos en los que estamos trabajando. Crear uno es algo que podemos hacer de dos maneras: en nuestro ordenador o en la nube\footnote{Decimos \emph{en la nube} para referirnos a que est\'a en internet, accesible a todos.}. Para hacer lo primero, vamos a abrir una Terminal y escribir \inlinecode{git init NombreDelProyecto}, en donde \textbf{NombreDelProyecto} es el nombre que nosotros querramos darle a la carpeta en donde colocaremos todos los archivos y documentos de nuestro proyecto. Luego entramos a la carpeta con \inlinecode{cd NombreDelProyecto} y aqu\'i podemos comenzar a trabajar.\\

La otra manera de comenzar un repositorio es accesando a una p\'agina como \href{https://github.com/}{GitHub}, \href{https://gitlab.com/}{GitLab}, \href{https://mercurial.selenic.com/}{Mercurial}, \href{https://bitbucket.org/}{BitBucket}, \href{http://sourceforge.net/}{SourceForge}, \href{http://www.codeplex.com/}{CodePlex}, etc. y creando un usuario en ella. Es conveniente mencionar que GitHub es la m\'as famosa, aunque de las anteriormente expuestas quiz\'a GitLab es la mejor para nuestros fines. Por eso, vamos a ir a GitLab, crearemos un usuario y seguiremos las instrucciones para crear un proyecto nuevo. Idealmente si nuestro proyecto tiene un nombre f\'acil y lo hacemos p\'ublico para que todos lo puedan ver y trabajar en \'el. Cuando nuestro repositorio ya est\'e creado, vamos a proveer quien nos est\'a dando el taller con nuestro usuario, y luego vamos a continuar.

\section{Clonando un repositorio}
Una vez logramos crear un repositorio en la nube, es momento de sincronizarlo con nuestro ordenador. Para ello vamos a ver que el proyecto que creamos en GitLab nos ofrece una direcci\'on \emph{HTTPS}. Vamos a copiar esta direcci\'on y vamos a escrbir lo siguiente en nuestra Terminal: \inlinecode{git clone} y luego vamos a pegar la direcci\'on que acabamos de copiar. Presionaremos enter y veremos como es que git crea una copia del proyecto en nuestro ordenador. Claro, como no hay nada en el proyecto a\'un, git nos advertir\'a de esto y solo crear\'a la carpeta.\\

Ahora que ya hemos comenzado, vamos a ver realmente para qu\'e es que nos sirve git. Para ello vamos a comenzar con un proyecto en particular y vamos a irlo desarrollando poco a poco entre todos los presentes.

\section{Proyecto}
El proyecto en el que vamos a trabajar ya est\'a creado en un repositorio de GitLab. Por esta raz\'on solo vamos a clonar el proyecto de esta direcci\'on: https://gitlab.com/zronyj/TQCA\_ej.git Al ya haber clonado el proyecto, entraremos en la carpeta correspondiente y comenzaremos a trabajar. Lo primero que haremos ser\'a entrar a la carpeta \textbf{Users} y crear un documento de texto que contenga nuestro nombre completo, nuestra fecha de cumplea\~nos y si somos estudiantes, catedr\'aticos o investigadores en una l\'inea diferente cada cosa. Luego vamos a guardar este documento bajo el siguiente nombre: \emph{apellidoNombre.txt} en donde el apellido y el nombre son \textbf{tu} apellido y \textbf{tu} nombre. Finalmente nos saldremos de la carpeta \textbf{Users} y quedaremos listos para el siguiente paso.\\

Crear funciones repartidas.

\section{Diferentes directorios de trabajo}
WD, Index (stage), HEAD

\section{Actualizando directorios: add \& commit}
add
commit

\section{Enviar cambios a la nube: push}
push

\section{Ramas}
checkout -b (crear)
checkout
checkout -d (borra)
checkout origin

\section{Actualizar y fusionar: pull \& merge}
pull
merge (rama)

\section{Revisar cambios: diff \& status}
diff
status

\section{Reemplazar cambios locales: checkout, fetch y \& reset}
checkout -- (archivo)

\section{Comentarios finales}

\newpage

\subsection{Glosario de comandos sencillos}
\begin{small}
\begin{itemize}
\item \textbf{git config}: configura git con tu nombre, tu direcci\'on de correo y tu editor de texto preferido.
\item \textbf{git init}: inicia un nuevo proyecto en git de manera local; un nuevo repositorio.
\item \textbf{git clone}: descarga los contenidos de un repositorio ya existente a tu ordenador y los deja listos para poder trabajar con ellos.
\item \textbf{git status}: revisa el estado actual de los archivos en el repositorio en comparaci\'on al \'ultimo commit.
\item \textbf{git diff}: muestra claramente las diferencias entre el estado actual de los archivos en el repositorio con respecto al \'ultimo commit. Muestra las diferencias l\'inea por l\'inea.
\item \textbf{.gitignore}: archivo que contiene el listado de todos los archivos del proyecto que deben ser ignorados por git.
\item \textbf{git commit}: guarda los cambios realizados a los archivos como una revisi\'on. Crea un registro y un \emph{id} que identifica a esa revisi\'on.
\item \textbf{git push}: empuja (sube) los commits locales a GitHub.
\item \textbf{git pull}: jala (descarga) los archivos del proyecto en GitHub, actualizando el estado de estos en nuestro ordenador.
\item \textbf{git log}: muestra toda la actividad de commits que ha habido en un repositorio/proyecto.
\item \textbf{git branch}: muestra a todas las ramas existentes o, si se le agrega un nombre despu\'es del comando, se crea una rama nueva en el proyecto.
\item \textbf{git merge}: une a una rama existente con otra rama que se debe de especificar como argumento del comando.
\item \textbf{git checkout}: si como argumento se agrega el nombre de una rama, se habr\'a trasladado a esa rama. Si como argumento se agrega el \emph{id} de un commit o la distancia de la cabeza a un commit con \inlinecode{HEAD$\sim$\#}, entonces nos situaremos en el commit seleccionado y podremos ver el estado del proyecto en ese momento.
\end{itemize}
\end{small}

\section*{Licencia}

\noindent \includegraphics{img/cc_big.png}

\noindent Taller de Qu\'imica Computacional Aplicada by \href{http://github.com/zronyj/TQCA}{Rony J. Letona} is licensed under a \href{http://creativecommons.org/licenses/by-sa/4.0/}{Creative Commons Attribution-ShareAlike 4.0 International License}.
Based on a work at \url{http://github.com/swcarpentry/bc}.

\end{document}
